\scanpage{99}
\begin{center}- \mspage{167} -\end{center}
\margin{IX}
$A$ \uline{lisse} sur $S$, et on désigne par
\[
(13.1.1)\qquad\underline E=\underline{\operatorname{Lie}}(A/S)
\]
le module localement libre sur $S$, image inverse par la section nulle du faisceau tangent relatif de $A$ sur $S$. Nous admettrons ici la définition de l'homomorphisme exponentiel
\[
(13.1.2)\qquad\exp:E\longrightarrow A\quad,
\]
où $E$ est le fibré vectoriel analytique sur $S$ dont le Module des sections holomorphes est $\underline E$. Cet homomorphisme se définit, au voisinage d'un point $s\in S$, d'abord au voisinage de la section nulle de $E$ par les méthodes bien connues en théorie des groupes formels en car. nulle, puis se prolonge de façon évidente par multiplicativité à $E$ tout entier. C'est un homomorphisme de groupes analytiques relatifs sur $S$, qui est \uline{étale} et a comme image le sous-espace ouvert $A^0$ de $A$ réunion des composantes neutres des fibres de $A$ sur $S$. Si $R$ est le sous-espace analytique noyau de $\exp$, on a donc une suite exacte de groupes analytiques relatifs :
\[
(13.1.3)\qquad
\diagram{\begin{tikzcd}0\arrow[r]&R\arrow[r]&E\arrow[r]&A^0\arrow[r]&0\end{tikzcd}}\quad,
\]
où le morphisme $E\longrightarrow A^0$ est surjectif et étale, donc $R$ est un groupe analytique relatif \uline{étale} sur $S$. C'est un sous-espace analytique de $E$, qui est fermé dans $E$ si et seulement si $A^0$ est séparé sur $S$, i.e. la section nulle de $A^0$ est fermée dans $A^0$.

\margin{485}
\margin{485}
\margin{98}
